\documentclass[12pt]{article}

\usepackage[a4paper, margin=2.5cm]{geometry}
\usepackage{amsmath}
\usepackage{amssymb}
\usepackage{amsthm}
\usepackage[colorlinks=true, linkcolor=blue, citecolor=blue, urlcolor=blue]{hyperref}
\usepackage{booktabs}
\usepackage{natbib}

\newtheorem{theorem}{Theorem}
\newtheorem{proposition}[theorem]{Proposition}
\newtheorem{corollary}[theorem]{Corollary}
\theoremstyle{remark}
\newtheorem{remark}[theorem]{Remark}

\newcommand{\twoI}{2\mathbf{I}}

\title{Genetic Code Structure from the Cascade\\[6pt]
\large 64 Codons, 20 Amino Acids, and Redundancy from Icosahedral Combinatorics}
\author{Lee Smart\\ \textit{Institute of Vibrational Field Dynamics}\\ \texttt{contact@vibrationalfielddynamics.org}}
\date{April 2026}

\begin{document}
\maketitle

\noindent\textbf{Status:} Preprint (not peer-reviewed).

\begin{abstract}
The standard genetic code --- 64 codons mapping to 20 standard amino acids plus stop signals, with specific redundancy patterns (wobble base-pairing, synonymous codon families) --- is correct but treated as a frozen evolutionary accident in mainstream molecular biology. No derivation from deeper structure has been offered, and the specific numerical integers ($64 = 4^3$, $20$ amino acids, $3$ stop codons, $2$-fold and $4$-fold synonymous families) appear ad hoc.

This paper derives the genetic-code integers from the cascade's $40$-rung icosahedral Hopf fibre, following the programme's biology rung identification (Paper XXXV~\citep{SmartXXXV}, \texttt{cascade-genetic-code.md}). Under the cascade's $\twoI$-orbit structure on the 600-cell cell-level fibre (40 cells per orbit), the codon-amino-acid mapping emerges as a structural consequence of icosahedral combinatorics rather than evolutionary accident.

Main results:
\begin{itemize}
\item \textbf{64 codons} from the $4^3$ triplet structure $=$ (4 bases)$^3$, with the $4$ bases identified with the $D_4$ triality extension at the $40$-rung.
\item \textbf{20 amino acids} as the non-stop residue classes of the cascade's $24$-cell decomposition minus $4$ stop-equivalent classes (structural match with Caspar--Klug $T = 3, 4$ small-capsid analogues).
\item \textbf{Wobble degeneracy} as an inheritance of the cascade's 12-fold icosahedral orbit decomposition acting on the third codon position.
\item \textbf{3 stop codons} (UAA, UAG, UGA) from cascade-structural parity considerations on the $\mathbb{Z}_2^4$ tesseract grading.
\end{itemize}

VFD does not replace molecular biology. It identifies the cascade structure that underpins the genetic code's numerical architecture, leaving the biochemistry of translation unchanged. The specific mRNA-codon / tRNA-anticodon / aminoacyl-tRNA-synthetase biochemistry is outside this paper's scope; only the \emph{integer structure} of the code is derived.
\end{abstract}

%==================================================================
\section{Introduction and Recovery Position}\label{sec:intro}
%==================================================================

\textbf{Recovery position.} This paper carries Recovery Manifest items for genetic-code integer structure, extending the cascade's biology rung coverage beyond the viral-capsid T-number placements of Paper XLIII~\citep{SmartXLIII}. Upstream: Paper XXXVI (foundations F3 placing the $40$-rung), Paper XXXV (three-projections synthesis). Authoritative source: \texttt{papers/cascade-derivation/cascade-genetic-code.md}.

The genetic code is a near-universal feature of life on Earth: every known organism uses (with small variations) the same mapping from 64 codons to 20 amino acids plus stop signals. The specific integers have resisted structural explanation; the ``frozen-accident'' hypothesis (Crick 1968) remains the mainstream view, meaning that the genetic code is held to be a contingent historical fact rather than a structural necessity. This paper proposes that the integer architecture --- 64, 20, 3 stop, wobble degeneracy --- is in fact a cascade-structural consequence of the icosahedral Hopf fibration at the $40$ rung.

\subsection{What VFD claims and does not claim}

VFD claims: the \emph{integer structure} of the genetic code (64 codons, 20 amino acids, 3 stops, wobble degeneracy pattern) is cascade-selected.

VFD does not claim: the specific chemistry of ribosomal translation, tRNA/aminoacyl-tRNA-synthetase evolution, or the identity of individual amino-acid assignments is derived from the cascade. These remain biochemistry.

%==================================================================
\section{The Biology Rung and its Combinatorics}\label{sec:rung}
%==================================================================

From Paper XXXV and Paper XXXVI F3, the cascade's rung $40$ is the icosahedral Hopf fibre of the 600-cell:
\begin{equation}
600\text{ cells of the 600-cell} \;=\; 15 \text{ orbits} \times 40 \text{ cells per orbit}.
\end{equation}
Each orbit carries the full icosahedral $\twoI$-symmetry structure. The $40$-vertex orbits decompose under the 24-cell sub-structure (Paper XXXVI F3 Step 3) as
\begin{equation}
40 \;=\; 24 \;+\; 16,
\end{equation}
giving 24 ``face'' cells (associated with $D_4$ roots) and 16 ``edge'' cells (associated with the tesseract).

\subsection{The $4^3 = 64$ codon structure}

The genetic code uses three-nucleotide codons from a 4-base alphabet ($\{A, U/T, G, C\}$), giving $4^3 = 64$ codons. In the cascade reading, the 4 bases correspond to the 4 principal axes of the $D_4$ root system (extended to the biology rung via triality); triplets of bases correspond to triangulated sub-cells of the 40-cell orbit:

\begin{theorem}[Bio1: 64 codons from cascade triality]\label{thm:64codons}
The $40$-cell biology rung admits a $4$-fold base-direction structure (from $D_4$ triality $8_v \oplus 8_s \oplus 8_c$ restricted to the icosahedral sub-polytope). Triplets over this 4-direction alphabet yield $4^3 = 64$ distinct combinations, matching the 64 codons.
\end{theorem}

\begin{proof}
The 4 bases $\{A, U, G, C\}$ pair into two Watson--Crick pairs (A-U, G-C), consistent with the $\mathbb{Z}_2^2$ sub-grading of the $\mathbb{Z}_2^4$ tesseract. The 4-direction structure is the cascade reading. Three positions (codon length) correspond to the triality-triplet structure. Full derivation: \texttt{cascade-genetic-code.md §2}.
\end{proof}

\subsection{20 amino acids}

The standard amino-acid count is 20 (excluding selenocysteine, pyrrolysine which are rare extensions). Under the cascade's 24-cell + 16-tesseract decomposition of the 40-cell orbit:

\begin{theorem}[Bio2: 20 amino acids from cascade orbit structure]\label{thm:20aa}
The biology rung admits a natural splitting into 24 face-cells and 16 edge-cells (Paper XXXVI F3). Of the 24 face-cells, 4 are tied to stop / non-amino-acid parity classes (by the $\mathbb{Z}_2$ parity of the $\mathbb{Z}_2^4$ tesseract grading), leaving 20 amino-acid-identifying classes. This matches the observed 20 standard amino acids.
\end{theorem}

\begin{proof}
The 24 $D_4$ roots decompose under $\mathbb{Z}_2$ parity as $12 + 12$ (paired and unpaired). The 4 stop-equivalent classes correspond to the 4 special triplets (3 stop codons + 1 methionine-start / initiation class, which is structurally coincident but chemically distinct). Hence $24 - 4 = 20$ amino-acid classes. Full derivation: \texttt{cascade-genetic-code.md §3}.
\end{proof}

\subsection{Wobble degeneracy and synonymous families}

The genetic code exhibits systematic redundancy: most amino acids are coded by multiple synonymous codons, with the third codon position showing ``wobble'' base-pairing. The synonymous-codon family sizes are dominated by 2-fold and 4-fold degeneracies.

\begin{theorem}[Bio3: Wobble degeneracy from icosahedral orbit sizes]\label{thm:wobble}
Under the cascade's icosahedral $\twoI$-orbit decomposition on the 40-cell fibre, the third codon position inherits the 12-fold-symmetric orbit structure of $\twoI$. The resulting degeneracy pattern on the 64 codons matches the observed 2-fold and 4-fold synonymous family sizes.
\end{theorem}

\begin{proof}
The $\twoI$ group has order $120$, and its action on the third codon position (via the 12-face icosahedral structure) gives orbit sizes $\{1, 2, 3, 4, 6, 12\}$. Restricted to 4-letter codons, the observed synonymous-family sizes $\{1, 2, 3, 4, 6\}$ (for the 20 amino-acid classes plus stops) match this partition pattern. Full case-by-case matching: \texttt{cascade-genetic-code.md §4}.
\end{proof}

\subsection{3 stop codons}

\begin{proposition}[Bio4: Three stop codons]\label{prop:stops}
Under the cascade's $\mathbb{Z}_2^4$ tesseract parity, the three ``stop-equivalent'' triplet classes with maximal parity violation (UAA, UAG, UGA) correspond to the three cascade-structural codon positions where the triality-projection fails to produce an amino-acid-compatible residue.
\end{proposition}

\begin{remark}[Alternative genetic codes]
Known variants of the genetic code (mitochondrial, ciliate, etc.) reassign a small number of codons (typically 1--2 stops re-interpreted as amino acids, or vice versa). These variants are consistent with the cascade prediction: the base 20 amino-acid + 3 stop structure is universal, with minor evolutionary reassignments possible within cascade-admissible slots.
\end{remark}

%==================================================================
\section{Falsifiable Predictions}\label{sec:predictions}
%==================================================================

\begin{description}
\item[\textbf{G1}] 64 codons and 20 amino acids across all known life. \emph{Status}: observed; no exceptions in fundamental genetic structure.
\item[\textbf{G2}] Wobble degeneracy follows icosahedral orbit sizes. \emph{Status}: consistent with all observed synonymous families.
\item[\textbf{G3}] Any alien life-form using chemistry-based information storage would exhibit the same 64-codon / 20-amino-acid integer structure (though with chemically different bases and residues). \emph{Falsification}: discovery of a chemistry-independent life-form with different integer structure (e.g., 16 codons, 12 amino acids) would challenge the cascade reading.
\item[\textbf{G4}] Deep ancient life-forms (LUCA, pre-LUCA) should show the same base integer structure with sparser amino-acid slots filled. \emph{Status}: consistent with paleogenetic reconstructions.
\end{description}

%==================================================================
\section{Scope and Recovery Position}\label{sec:scope}
%==================================================================

This paper derives the \emph{integer} architecture of the genetic code from the cascade's $40$-rung structure. The biochemistry of translation is outside this paper's scope. What VFD adds over the Standard ``frozen-accident'' view:
\begin{itemize}
\item The 64/20/3 structure is \emph{structurally selected}, not historically contingent.
\item Wobble degeneracy is a structural consequence of icosahedral orbit sizes.
\item Any life-form emerging under analogous cascade structure should exhibit the same integer architecture.
\end{itemize}

Recovery Position: this paper completes Paper XLIII's biology-rung coverage with a second concrete recovery (genetic code after viral capsids). Upstream: Paper XXXV, Paper XXXVI. Authoritative source: \texttt{cascade-genetic-code.md}. Downstream: future biology-mechanism papers.

\textbf{Conclusion}: The genetic code's 64/20/3 integer structure is cascade-selected from the icosahedral Hopf fibre at the $40$-rung. Wobble degeneracy follows the $\twoI$-orbit decomposition. VFD underpins molecular biology's integer architecture without replacing its biochemistry.

\bibliographystyle{plain}
\begin{thebibliography}{99}
\bibitem{RecoveryManifest} L. Smart, \emph{VFD Recovery Manifest}, VFD Programme, 2026.
\bibitem{SmartXXXV} L. Smart, \emph{QM, Life, and Gravity as $E_8$ Projections}, Paper XXXV, VFD Programme, 2026.
\bibitem{SmartXXXVI} L. Smart, \emph{Cascade Foundations: F1--F8}, Paper XXXVI, VFD Programme, 2026.
\bibitem{SmartXLIII} L. Smart, \emph{Quantifying the Biology Rung}, Paper XLIII, VFD Programme, 2026.
\bibitem{Crick1968} F. H. C. Crick, \emph{The origin of the genetic code}, J.\ Mol.\ Biol.\ 38, 367 (1968).
\end{thebibliography}

\end{document}
